\section{CFD setup: differences from REPORT-4}
\label{app:setup}

The runs follow the study-A workflow of REPORT-4 (geometry, mesh, models, boundary conditions and
solvers). Only these points differ.

\begin{itemize}
\item \textbf{Domain}: built in Ansys Discovery from the 3D STEP of each nose: the end disc is pulled to
$x = \SI{2520}{mm}$, a $6120 \times \SI{3600}{mm}$ rectangle on the XY plane of the origin is subtracted from
it. Area \SI{21.8233}{m^2} (A) and \SI{21.8269}{m^2} (B), within \SI{0.001}{\percent} of the domains
of \file{CFD/make_domain_2d.py}.
\item \textbf{Named selections}: the body line is split at the end of the coupler, $x = \SI{720}{mm}$,
into \zone{wall\_coupler} and \zone{wall\_body}; seven zones in all.
\item \textbf{Mesh}: \zone{wall\_body} edge sizing hard, from \SI{0.55}{mm} at $x = \SI{720}{mm}$ to
\SI{1.65}{mm} at the outlet; inflation scoped on the face with the three walls as boundary.
\item \textbf{Fluent}: 2024~R2 started on its own and set up by script (\file{CFD/study720/setup_fluent.py},
\file{run_queue.py}, PyFluent). Turbulence intensity set to \SI{1}{\percent} (the standalone default is
\SI{5}{\percent}). Report definitions \texttt{cd-total} (three walls) and \texttt{cd-nose}
(\zone{wall\_nose}); residual stopping off.
\item \textbf{Iterations}: 500 first order at Courant 1, 500 at Courant 5, then second order at Courant 5
up to 4000. $M = 0.95$: pressure-based coupled solver, 1200 iterations. \SI{300}{K} walls: 3500 more
iterations from the adiabatic solution.
\end{itemize}

Scripts, wall exports, $C_D$ histories, figures and meshes are in \file{CFD/study720/} of the repository
(README there); the case and data files stay on the DLL PC.
